
\chapter{Methodology}

\startcontents[chapters]

\section{Agreement sources}
The corpus used for this study was the 300M subcorpus of the National Corpus of Polish \citep{przepiorkowski_narodowy_2012}, which was parsed using Trankit \citep{trankit}. The dependency trees (conforming to the Universal Dependency standard, henceforth UD) obtained this way were searched for words, which could trigger hybrid agreement: profession nouns, socially gendered nouns and depreciative forms of M1 nouns. The subsections below describe how those nouns were chosen.

\subsection{Profession nouns}
Nouns in this category came from two studies about feminitives in Polish. The first was a survey reported by \cite{szpyra-kozlowska_premiera_2019}, where participants were asked to create feminitives from 50 masculine nouns, that in modern Polish do not have feminine counterparts at all or those counterparts are rarely used. For each noun participants had the option to give no answer. The second study was an experiment conducted by \cite{szpyra-kozlowska_czy_2022}, where they tested whether certain feminitives are less accessible by asking participants to try to pronounce them. They chose 20 feminitives, which are formed from masculine nouns that end in consonant clusters such as /kt/, /tr/ or /rg/. According to \cite{klemensiewicz_tytuly_1957} adding the -ka suffix, the function of which can be to create the feminitve, makes the word difficult to pronounce. Those two studies give a list of 60 nouns in total, for which it is said to be difficult to form feminine counterparts. It is therefore likely that, if there was a female referent of a noun from this list, the masculine noun would be used, leading to hybrid agreement. 

Both of the studies determine in some way which of the tested nouns have less available feminitives than others: in \cite{szpyra-kozlowska_czy_2022} those are the nouns for which there are the highest error rates in pronunciation, whereas in \cite{szpyra-kozlowska_premiera_2019} -- the nouns for which the most people refused to propose a feminitive. Taking this approach has, however, two downsides: first, it requires an arbitrary threshold that separates the nouns that have an easily formed feminitive from the nouns for which creating the feminine counterpart is more difficult. Second, this approach significantly limits the number of nouns for analysis. Considering this, the current study used all of the nouns from the two studies on Polish feminitves.

For ease of reference, all of those nouns are called \textsl{profession} nouns, despite not all of them actually referring to professions. Meanings of all nouns are presented in Table \ref{tab:noun-translations} in the Appendix.

\subsection{Socially gendered nouns}
The nouns for this category were found using the Polish Academy of Sciences Great Dictionary of Polish available online.\footnote{https://wsjp.pl/} The advanced search option allowed to search for nouns in the thematic category \textsc{human as a physical being: human physicality: sex}. This search yielded 586 results, which then had to be manually filtered to find the nouns that only refer to humans (and not, for example, abstract nouns or nouns referring to body parts) and that had incongruent grammatical and semantic gender. The 19 nouns selected this way can be found in Table \ref{tab:noun-translations} in the Appendix.

\subsection{Depreciative forms}
For this category no list of nouns was created, since the corpus was annotated and it was possible to look up the depreciative forms by checking the words' morphological description. It is possible to create depreciative forms of surnames, however these were excluded from the analysis, because there was no way of determining the referents and therefore no way of ensuring that there is hybrid agreement. Surname \textsl{Machała} found in the corpus can serve as an example: let us say that the \textsl{Machała} family bought a new car. Assuming the family consists of at least one man, the speaker can use both (\ref{ex:Machałowie}) and (\ref{ex:Machały}) to convey that information, depending on their attitude they have towards the situation: (\ref{ex:Machałowie}) if the speaker is neutral or positive and (\ref{ex:Machały}) if they are negative. However, if there are only women in the family, the speaker can only use (\ref{ex:Machały}). Without the context it is unclear whether \textsl{Machały} in (\ref{ex:Machały}) is used as a depreciative or simply to refer to women and therefore we cannot decide if this is an example of hybrid agreement.

\begin{exe}
  \ex
  \begin{xlist}
    \ex\label{ex:Machałowie}
    \gll Machał-owie kupi-li nowe auto.\\
    Machała-pl.nom.m1.ndepr buy-3pl.m1.past new car\\
    \ex\label{ex:Machały}
    \gll Machał-y kupi-ły nowe auto.\\
    Machała-pl.nom.m1.depr buy-3pl.f.past new car\\
    \trans ``Machałas bought a new car.''
  \end{xlist}
\end{exe}

Morphological analyzer Morfeusz \citep{kieras_morfeusz_2017} can mark forms as surnames which made it easy to exclude them from the analysis. 

\subsection{Finding agreement sources in the corpus}
The code used for extracting corpus data in this study can be found in a GitHub repository.\footnote{https://github.com/bmagdab/hybrid-agreement} The general idea of how the code works is presented in this chapter.

After choosing the nouns to be searched in the corpus, it was necessary to find all of their forms and the forms' morphological descriptions, which was done using the morphological analyzer Morfeusz. First the corpus was searched for occurrences of nouns, which could be sources of hybrid agreement. Those nouns were saved if they were in the list of gendered or profession noun forms or if they were marked as depreciative in their morphological description. To ensure that no homonyms were accidentaly saved for later analysis, the morphological description of a found word was compared with the one given by Morfeusz. This way, for example, the word \textsl{premiera} was excluded if it was an instance of the lexeme \textsl{premiera}, meaning \textsl{premiere} but not if it was an instance of the lexeme \textsl{premier}, meaning \textsl{prime minister}, in the singular accusative form.

Once it was sure that the right noun was found, the dependency tree was searched for words that can agree their gender with that of the noun. This was the head of the noun (which could be an example of predicative agreement), adjectives attached to the noun (which could be examples of attributive agreement), and relative pronouns referring to the noun. Based on the UD annotation guidelines, all dependents of the word were searched for a dependency label with a subtype \textsc{:relcl}, which in the UD annotation scheme mark the head of the relative clause. The relative pronoun was attached to the head of the relative clause and had a UD feature \textsc{PronType=Rel}.\footnote{https://universaldependencies.org/workgroups/newdoc/relative\_clauses.html} Once the possible agreement targets were found, the data could be used to look for occurrences of gender agreement.

\section{Listing agreement occurrences}
To analyse hybrid gender agreement in the found examples, it is necessary to mark the alternative to the grammatical gender of each agreement source. For all profession nouns that alternative gender was feminine, because the only chance of hybrid agreement for those noun is when the gender of the referent is feminine. For socially gendered nouns the alternative gender was assigned depending on the lexeme.

Once the alternative gender was determined, it was necessary to exclude plural nouns for which the gender discrepancy was between feminine and neuter. This is because for all types of agreement tested here, the plural feminine and plural neuter forms of targets of agreement are syncretic, therefore there is no hybrid agreement in those cases. This happens, for instance, in (\ref{ex:plural-feminine-neuter}), where in singular ((\ref{ex:sg.n}) and (\ref{ex:sg.f})) there is a difference between the feminine and neuter forms of the word \textsl{młoda}. In plural (\ref{ex:pl}), however, the form is the same for both neuter and feminine, so hybrid agreement cannot be observed. Examples like this were therefore excluded.

\begin{exe}
  \ex\label{ex:plural-feminine-neuter}
  \begin{xlist}
    \ex\label{ex:sg.n}
    \gll młod-e dziewczę\\
    young-sg.n.nom girl.sg.n.nom\\
    \ex\label{ex:sg.f}
    \gll młod-a dziewczę\\
    young-sg.f.nom girl.sg.n.nom\\
    \ex\label{ex:pl}
    \gll młod-e dziewczę-ta\\
    young-pl.nom girl-pl.n.nom\\
    \trans ``young girl''
  \end{xlist}
\end{exe}

The next step was to find occurrences of different types of agreement. The predicative agreement of a given noun was found in its head on three conditions: first, dependency label between the noun and the head had to be \textsc{nsubj}, which is the UD label for nominal subjects. Second, the head had to be a verb that inflects for gender, which is only true for verbs annotated as past participles,\footnote{The annotation scheme used for NKJP uses the past participle label not only for past participles, but also for parts of past tense and conditional mood forms \citep{przepiorkowski_narodowy_2012}.} active and passive participles. Third, the noun has to be in the nominative case. This condition helps filter out situations where a numeral phrase is in subject position, because then the verb does not agree with the subject -- it is in 3rd person, singular and neuter form, regardless of the noun. Therefore, predicative agreement cannot be observed in those situations at all and they should be excluded.

The search was simpler with attributive and relative pronoun agreement. The only requirement for attributive agreement was that the word is an adjective, which was verified by looking for the tag \textsc{adj} in the morphological description of the dependent. As for relative pronouns, all of the ones found earlier were included in the analysis, unless they did not inflect for gender (which was the case, for example, for the word \textsl{gdzie}, meaning \textsl{where}). 

The sentence in (\ref{ex:agreement-extraction}) serves as an example of the process of extracting agreement described in this chapter. The word, which could trigger hybrid agreement (but here does not, since the referent is male) is \textsl{premierem}. There are two words in the sentence that agree with this word in terms of gender: \textsl{jedynym} and \textsl{który}. All of the information saved for each instance of agreement is shown below. 

% NKJP_1M: 3571
\begin{exe}
  \ex\label{ex:agreement-extraction}
  \gll Nie jest prawd-ą, że premier Jerzy Buzek jest jedyn-y-m- premier-em, któr-y wystąpi-ł w Senac-ie.\\
  not be.3sg.pres truth.f.sg-inst that prime.minister.sg.m1.nom Jerzy Buzek be.3sg.pres only-m1-inst-sg prime.minister.m1-sg.inst which-sg.nom.m1 appear-3sg.m1.past in senate.m2-sg.loc \\
  \trans ``It's not true, that prime minister Jerzy Buzek is the only one that has appeared in the Senate.''
\end{exe}

\begin{tblr}{
    hlines, vlines,
    cell{2}{1} = {r=5,c=1}{c},
    cell{7}{1} = {r=3,c=1}{c},
    cell{1,10}{1} = {c=2}{l},
    cell{2-6}{3} = {c=2}{l},
  }
  & & Occurrence 1 & Occurrence 2 \\
  source of agreement & form & premierem & \\
  & lexeme & premier & \\
  & morphological description & subst:sg:inst:m1 & \\
  & grammatical gender & m1 & \\
  & other gender & f & \\
  agreement target & form & jedynym & który \\
  & morphological description & adj:sg:inst:m1:pos & adj:sg:nom:m1:pos \\
  & gender & m1 & m1 \\
  type of agreement & & attributive & relative pronoun\\
\end{tblr}
